\documentclass[a4paper,12pt]{article}
\usepackage[T2A]{fontenc}
\usepackage[utf8]{inputenc}
\usepackage[russian]{babel}
\usepackage{amsmath,amssymb,mathtools}
\usepackage{icomma}
\usepackage[a4paper,top=25mm,bottom=24mm,left=26mm,right=25mm,headheight=24pt]{geometry}
\usepackage{microtype}
\usepackage{tikz}
\usetikzlibrary{calc,arrows.meta,positioning,shapes.geometric}
\usepackage{booktabs,tabularx}
\usepackage{needspace}
\usepackage{enumitem}
\usepackage{fancyhdr}
\usepackage{setspace}
\usepackage{xcolor}
\usepackage[hidelinks]{hyperref}
\definecolor{accent}{HTML}{008080}
\definecolor{darkaccent}{HTML}{005858}
\definecolor{textgray}{HTML}{343C42}
\definecolor{softline}{HTML}{B9D7D5}
\setlength{\parindent}{0pt}
\setlength{\parskip}{5pt}
\onehalfspacing
\setlist[itemize]{leftmargin=6mm,itemsep=1pt,topsep=2pt}
\setlist[enumerate]{leftmargin=7mm,itemsep=3pt,topsep=3pt}
\pagestyle{fancy}
\fancyhf{}
\fancyhead[L]{\footnotesize\color{textgray}Проценты и распределительный закон}
\fancyhead[R]{\footnotesize\color{textgray}Надя Клименко}
\renewcommand{\headrulewidth}{0.35pt}
\renewcommand{\headrule}{\hbox to\headwidth{\color{softline}\leaders\hrule height \headrulewidth\hfill}}
\fancyfoot[C]{\footnotesize\color{textgray}\thepage}
\newcommand{\sect}[1]{\Needspace{8\baselineskip}\section*{\color{darkaccent}#1}\vspace{-2mm}\noindent{\color{softline}\rule{\linewidth}{0.5pt}}\vspace{1mm}}
\newcommand{\markkey}[1]{\textcolor{darkaccent}{\textbf{#1}}}
\begin{document}
\begin{titlepage}
\thispagestyle{empty}
\centering
\vspace*{21mm}
{\Large\color{darkaccent} МАТЕМАТИКА\par}
\vspace{20mm}
{\LARGE\color{textgray} Чек-лист\par}
\vspace{12mm}
{\LARGE\bfseries\color{accent} Проценты, распределительный закон\\[3mm] и подобные слагаемые\par}
\vspace{13mm}
{\large Повторение ключевых приёмов занятия\par}
\vfill
{\large Надя Клименко\par}
\vspace{10mm}
{\normalsize 14 ноября 2026 г.\par}
\vspace{17mm}
{\small Лёвин Артём Александрович\par}
{\small эксклюзивно для Нади Клименко\par}
\vspace{16mm}
\begin{tikzpicture}[x=1cm,y=1cm]
  \draw[accent,line width=0.9pt] (0,0) .. controls (1.0,0.45) and (2.0,-0.38) .. (3,0)
    .. controls (4.1,0.42) and (5.1,-0.35) .. (6.2,0);
\end{tikzpicture}
\vspace{3mm}
\end{titlepage}
\setcounter{page}{1}

\sect{1. На сколько процентов число стало больше}
Если число изменилось от \(b\) до \(a\), где \(a>b>0\), процент увеличения рассчитывается \textbf{от первоначального числа} \(b\):
\[
\boxed{\text{Увеличение, \%}=\frac{a-b}{b}\cdot100\%.}
\]
Для цены запоминайте: \(\displaystyle\frac{\text{стало}-\text{было}}{\text{было}}\cdot100\%\). В знаменателе всегда находится исходная величина.

\markkey{Пример.} Цена выросла с 240 до 300 рублей. Тогда
\[
\frac{300-240}{240}\cdot100\%=\frac{60}{240}\cdot100\%=25\%.
\]
\markkey{Проверка.} 25\% от 240 рублей --- это 60 рублей; \(240+60=300\).

\sect{2. Распределительный закон}
\markkey{Формулы.} Число перед скобками умножается на \emph{каждое} слагаемое или вычитаемое:
\[
\boxed{a(b+c)=ab+ac,\qquad a(b-c)=ab-ac.}
\]
Перед вычислением посмотрите, получится ли удобно сократить дроби после умножения. Если да, раскройте скобки; если нет, иногда проще сначала выполнить действие внутри них.

\markkey{Пример.} Вычислим без сложения дробей в скобках:
\[
\left(\frac49+\frac7{36}\right)\cdot36
=\frac49\cdot36+\frac7{36}\cdot36
=16+7=\boxed{23}.
\]
\markkey{Важно.} Делить и сокращать можно \emph{множители} в произведении. Нельзя сокращать отдельные слагаемые через знак плюс.

\sect{3. Умножение смешанных чисел без лишних действий}
Смешанное число --- сумма целой и дробной частей: \(3\frac19=3+\frac19\).
\[
3\frac19\cdot5
=\left(3+\frac19\right)\cdot5
=15+\frac59
=\boxed{15\frac59}.
\]
Ещё один пример: \(\displaystyle5\frac15\cdot5=(5+\frac15)\cdot5=25+1=\boxed{26}.\)

Скобки могут содержать несколько смешанных чисел:
\[
\begin{aligned}
\left(2\frac34+7\frac13\right)\cdot6
&=\left(2+\frac34+7+\frac13\right)\cdot6\\
&=12+\frac{18}{4}+42+2
=\boxed{60\frac12}.
\end{aligned}
\]
Преимущество приёма: можно избежать перевода всех смешанных чисел в неправильные дроби и поиска общего знаменателя.

\sect{4. Зачем нужны скобки при вычитании}
\markkey{Правило «ВУПС».} \textbf{В}ычитаемое, \textbf{У}множаемое или \textbf{П}одставляемое выражение из нескольких слагаемых берите в \textbf{С}кобки.

При разложении смешанного числа, которое вычитают, скобки обязательны:
\[
10-2\frac1{11}=10-\left(2+\frac1{11}\right).
\]
\markkey{Пример.} Используем скобки и распределительный закон:
\[
\begin{aligned}
\left(10-2\frac1{11}\right)\cdot11
&=\left(10-\left(2+\frac1{11}\right)\right)\cdot11\\
&=110-22-1=\boxed{87}.
\end{aligned}
\]
\markkey{Типичная ошибка.} Запись \(10-2+\frac1{11}\) меняет значение исходного выражения. Минус перед скобками относится ко \emph{всей} сумме: \(a-(b+c)=a-b-c\).

\sect{5. Распределительный закон в обратную сторону}
Если один и тот же множитель встречается в нескольких слагаемых, вынесите его за скобки:
\[
\boxed{ab+ac=a(b+c),\qquad ab-ac=a(b-c).}
\]
\markkey{Пример.} Повторяется смешанное число \(7\frac4{13}\):
\[
\begin{aligned}
7\frac4{13}\cdot5\frac57+7\frac4{13}\cdot7\frac27
&=7\frac4{13}\cdot\left(5\frac57+7\frac27\right)\\
&=7\frac4{13}\cdot13
=\frac{95}{13}\cdot13=\boxed{95}.
\end{aligned}
\]
\markkey{Смотрите вперёд.} Быстрое сложение в скобках часто делает дальнейшее умножение значительно проще.

\sect{6. Приведение подобных слагаемых}
\textbf{Подобные слагаемые} имеют одинаковую буквенную часть. Складывайте или вычитайте их \emph{коэффициенты}, сохраняя буквенную часть.
\[
\boxed{\alpha x+\beta x=(\alpha+\beta)x.}
\]
\markkey{Примеры.}
\[
\begin{aligned}
\frac37a+\frac27a&=\left(\frac37+\frac27\right)a=\boxed{\frac57a};\\[2mm]
\frac4{11}a+a&=\left(\frac4{11}+1\right)a=\boxed{\frac{15}{11}a};\\[2mm]
\frac35x+\frac2{15}x-\frac4{15}x
&=\frac{9+2-4}{15}x=\boxed{\frac7{15}x}.
\end{aligned}
\]
У одиночной буквы коэффициент равен единице: \(a=1\cdot a\). В ответе обычно записывают сначала коэффициент, затем букву.

Если перед скобкой стоит минус, сохраняйте её до правильного раскрытия:
\[
\begin{aligned}
\frac9{14}c-\left(\frac3{14}c+\frac27c\right)
&=\frac9{14}c-\left(\frac3{14}c+\frac4{14}c\right)\\
&=\left(\frac9{14}-\frac7{14}\right)c
=\boxed{\frac17c}.
\end{aligned}
\]

\sect{7. Распределительный закон в уравнениях}
\markkey{Алгоритм.} Раскройте скобки, приведите подобные слагаемые, перенесите известные числа в одну сторону и найдите неизвестное. В конце подставьте результат в первоначальное уравнение.

\markkey{Пример.}
\[
\begin{aligned}
\frac34(20-16y)&=3,\\
15-12y&=3,\\
12y&=15-3=12,\\
\boxed{y=1}.&
\end{aligned}
\]
\markkey{Проверка.} \(\frac34(20-16\cdot1)=\frac34\cdot4=3\). Равенство верное.

\sect{8. Мини-памятка перед ответом}
\begin{itemize}
\item При процентном увеличении делите разницу на значение \textbf{«было»}.
\item При раскрытии скобок умножайте на \textbf{каждое} слагаемое.
\item При вычитании суммы сохраняйте \textbf{скобки} и правильно меняйте знаки.
\item При повторяющемся множителе ищите возможность \textbf{вынести его за скобки}.
\item При приведении подобных складывайте \textbf{коэффициенты}; буква остаётся прежней.
\item Проверяйте, что знаменатели дробей отличны от нуля; в уравнении выполняйте подстановку.
\end{itemize}

\clearpage
\thispagestyle{fancy}
\begin{center}
{\Large\bfseries\color{darkaccent}Алгоритм: как выбрать удобное преобразование}
\end{center}
\vspace{13mm}
\begin{center}
\begin{tikzpicture}[
  every node/.style={align=center,font=\small},
  action/.style={draw=softline,very thick,rounded corners=2pt,text width=94mm,minimum height=13mm,fill=white,inner sep=3pt},
  choice/.style={draw=accent,very thick,diamond,aspect=3.6,text width=70mm,inner sep=1.5pt,minimum height=17mm},
  endpoint/.style={draw=accent,very thick,rounded corners=13pt,text width=77mm,minimum height=11mm,inner sep=3pt},
  arr/.style={-{Latex[length=2.4mm]},color=darkaccent,line width=0.9pt}
]
\node[endpoint] (begin) at (0,0) {Рассмотрите выражение};
\node[choice] (one) at (0,-2.35) {Есть число, умножающее сумму или разность?};
\node[action] (expand) at (0,-4.70) {Раскройте скобки, если это упрощает расчёт};
\node[choice] (two) at (0,-7.15) {Повторяется общий множитель?};
\node[action] (factor) at (0,-9.5) {Вынесите общий множитель за скобки};
\node[action] (terms) at (0,-11.55) {Соберите подобные слагаемые и вычислите коэффициенты};
\node[endpoint] (finish) at (0,-13.6) {Получите результат и проверьте его};
\draw[arr] (begin) -- (one);
\draw[arr] (one) -- node[right,xshift=2mm]{да} (expand);
\draw[arr] (expand) -- (two);
\draw[arr] (two) -- node[right,xshift=2mm]{да} (factor);
\draw[arr] (factor) -- (terms);
\draw[arr] (terms) -- (finish);
\draw[arr] (one.east) -- ++(2.1,0) node[above,xshift=-10mm]{нет} |- (two.east);
\draw[arr] (two.west) -- ++(-2.2,0) node[above,xshift=10mm]{нет} |- (terms.west);
\end{tikzpicture}
\end{center}
\clearpage

\sect{Тренировочный блок --- 10 заданий}
\textit{Путь А.} Решайте от простого к сложному. Для вычислений выбирайте приём, который сокращает число действий. В заданиях на упрощение записывайте коэффициент перед буквой.
\begin{enumerate}[label=\textbf{\arabic*.},itemsep=5pt]
\item На сколько процентов число 600 больше числа 480?
\item Цена товара выросла с 320 до 400 рублей. На сколько процентов увеличилась цена?
\item Вычислите: \(\displaystyle\left(\frac56+\frac7{18}\right)\cdot18\).
\item Вычислите, используя разложение смешанного числа: \(\displaystyle4\frac17\cdot7\).
\item Вычислите удобным способом: \(\displaystyle\left(1\frac34+2\frac16\right)\cdot12\).
\item Вычислите, сохраняя необходимые скобки: \(\displaystyle\left(12-3\frac18\right)\cdot8\).
\item Вычислите с вынесением общего множителя:
\[
6\frac38\cdot2\frac25+6\frac38\cdot5\frac35.
\]
\item Приведите подобные: \(\displaystyle\frac5{12}a+\frac7{18}a-\frac19a\).
\item Упростите выражение: \(\displaystyle\frac{11}{15}b-\left(\frac25b+\frac1{15}b\right)\).
\item Решите уравнение и проверьте ответ: \(\displaystyle\frac56(18-6x)=5\).
\end{enumerate}
\vfill
{\color{darkaccent}\small\textbf{Самопроверка:} в заданиях 1--2 укажите, от какой величины находили процент; в заданиях 6 и 9 внимательно следите за скобками; в задании 10 подставьте корень в исходное уравнение.}
\clearpage

\sect{Ответы к тренировочному блоку}
\renewcommand{\arraystretch}{1.6}
\begin{center}
\begin{tabularx}{0.82\linewidth}{@{}>{\bfseries}c X@{}}
\toprule
№ & Ответ \\
\midrule
1 & \(25\%\) \\
2 & \(25\%\) \\
3 & \(22\) \\
4 & \(29\) \\
5 & \(47\) \\
6 & \(71\) \\
7 & \(51\) \\
8 & \(\dfrac{25}{36}a\) \\
9 & \(\dfrac4{15}b\) \\
10 & \(x=2\) \\
\bottomrule
\end{tabularx}
\end{center}
\vspace{10mm}
\sect{Как проверить себя}
\begin{itemize}
\item \textbf{Проценты:} результат умножьте на исходное число и убедитесь, что получили разницу.
\item \textbf{Дроби:} попробуйте выполнить расчёт другим способом; ответы должны совпасть.
\item \textbf{Буквы:} проверьте отдельно коэффициент и правильность знаков.
\item \textbf{Уравнение:} подставьте найденное значение в обе части.
\end{itemize}
\vfill
\begin{center}
{\color{textgray}\footnotesize Лёвин Артём Александрович\quad\textbullet\quad эксклюзивно для Нади Клименко}
\end{center}
\end{document}
